% Chapter 03: Validation Rules
% Writing semantic validators beyond JSON Schema

\chapter{Validation Rules}
\label{ch:validation-rules}

% ============================================================
\section{Two Classes of Validation}
\label{sec:validation-classes}
% ============================================================

Validation in the Domain Kit model operates at two levels:

\begin{description}
  \item[Structural validation (schema)] is performed by JSON Schema at parse
    time. It enforces field presence, types, formats, and per-step correctness.
    This is covered in Chapter~\ref{ch:authoring-schemas}.

  \item[Semantic validation (business rules)] is performed by Kit-specific
    validators at parse time or plan time. These are executable binaries that
    enforce domain-specific invariants that cannot be expressed in JSON Schema
    alone.
\end{description}

\paragraph{When to use semantic validators:}

\begin{itemize}
  \item \textbf{Cross-step invariants:} ``Every provision step must be followed
    by a cleanup step.''
  \item \textbf{Contextual constraints:} ``If \texttt{timeout} exceeds SLA,
    emit a warning.''
  \item \textbf{Domain-specific business rules:} ``Every runbook in the
    compliance Kit must declare an owner in \texttt{meta.roles}.''
  \item \textbf{Complex conditions:} ``If step A sets variable X, step B (if
    present) must reference X.''
\end{itemize}

If the rule can be expressed in JSON Schema, use JSON Schema. Semantic
validators are for everything else.

% ============================================================
\section{Writing Semantic Validators}
\label{sec:writing-validators}
% ============================================================

A semantic validator is a standalone executable that reads the parsed runbook
(or the lowered runbook) on stdin and writes a list of validation errors and
warnings to stdout.

\paragraph{Input contract:}

The validator receives the runbook as YAML on stdin. For parse-time validators,
this is the Kit-authored runbook. For plan-time validators, this is the lowered
\texttt{runbook/v1} document.

\paragraph{Output contract:}

The validator writes a JSON array of validation messages to stdout. Each message
has:

\begin{description}
  \item[\texttt{file}] --- File path (or \texttt{"<stdin>"} if from a pipe).
  \item[\texttt{line}] --- Line number where the issue occurs (1-indexed).
  \item[\texttt{column}] --- Column number (1-indexed, optional).
  \item[\texttt{severity}] --- \texttt{"error"} or \texttt{"warning"}.
  \item[\texttt{message}] --- Human-readable description of the issue.
  \item[\texttt{rule}] --- Validator rule identifier (e.g.,
    \texttt{"compliance/owner-required"}).
\end{description}

\paragraph{Example output:}

\begin{minted}{json}
[
  {
    "file": "runbook.yaml",
    "line": 3,
    "column": 1,
    "severity": "error",
    "message": "Compliance runbooks must declare an owner in meta.roles.owner",
    "rule": "compliance/owner-required"
  },
  {
    "file": "runbook.yaml",
    "line": 15,
    "severity": "warning",
    "message": "Timeout '4h' exceeds recommended SLA of 2h",
    "rule": "compliance/timeout-sla"
  }
]
\end{minted}

If validation passes, the validator writes an empty array \texttt{[]} to stdout.

\paragraph{Exit code contract:}

\begin{description}
  \item[\texttt{0}] --- Validation passed (no errors).
  \item[\texttt{1}] --- Validation failed (at least one error).
  \item[\texttt{2+}] --- Internal validator error (e.g., malformed input, panic).
\end{description}

Warnings do not cause a non-zero exit code. Only errors cause exit code 1.

% ============================================================
\section{Parse-Time Validators}
\label{sec:parse-validators}
% ============================================================

Parse-time validators run when \texttt{gert validate} is called, \emph{before}
lowering. They see the Kit-authored runbook and can enforce authoring-layer
invariants.

\paragraph{Use cases:}

\begin{itemize}
  \item Enforce that required metadata fields are present (e.g.,
    \texttt{meta.roles.owner}).
  \item Validate that step IDs follow naming conventions (e.g., no spaces,
    kebab-case).
  \item Check that Kit-specific fields are used correctly (e.g.,
    \texttt{approvers} list is not empty).
\end{itemize}

\paragraph{Example: \texttt{owner-required} validator (pseudocode):}

\begin{minted}{go}
package main

import (
  "encoding/json"
  "fmt"
  "gopkg.in/yaml.v3"
  "os"
)

type ValidationMessage struct {
  File     string `json:"file"`
  Line     int    `json:"line"`
  Severity string `json:"severity"`
  Message  string `json:"message"`
  Rule     string `json:"rule"`
}

func main() {
  var runbook map[string]interface{}
  yaml.NewDecoder(os.Stdin).Decode(&runbook)

  var messages []ValidationMessage

  meta, ok := runbook["meta"].(map[string]interface{})
  if !ok || meta["roles"] == nil {
    messages = append(messages, ValidationMessage{
      File:     "<stdin>",
      Line:     1,
      Severity: "error",
      Message:  "Compliance runbooks must declare meta.roles.owner",
      Rule:     "compliance/owner-required",
    })
  } else {
    roles := meta["roles"].(map[string]interface{})
    if roles["owner"] == nil {
      messages = append(messages, ValidationMessage{
        File:     "<stdin>",
        Line:     3,
        Severity: "error",
        Message:  "meta.roles.owner is required for compliance runbooks",
        Rule:     "compliance/owner-required",
      })
    }
  }

  json.NewEncoder(os.Stdout).Encode(messages)
  if len(messages) > 0 && hasError(messages) {
    os.Exit(1)
  }
}

func hasError(messages []ValidationMessage) bool {
  for _, m := range messages {
    if m.Severity == "error" {
      return true
    }
  }
  return false
}
\end{minted}

% ============================================================
\section{Plan-Time Validators}
\label{sec:plan-validators}
% ============================================================

Plan-time validators run \emph{after} lowering, when the planner is about to
produce an ExecutionPlan. They see the lowered \texttt{runbook/v1} document
(core primitives only) and can enforce cross-step invariants.

\paragraph{Use cases:}

\begin{itemize}
  \item Enforce step ordering rules (e.g., ``cleanup must follow provision'').
  \item Validate that lowered steps have required governance fields populated.
  \item Check that synthesized variables are referenced correctly.
\end{itemize}

\paragraph{Example: \texttt{cleanup-follows-provision} validator:}

This validator checks that for every step with ID matching \texttt{provision-*},
there exists a subsequent step with ID matching \texttt{cleanup-*}.

\begin{minted}{go}
func validateCleanupFollowsProvision(runbook map[string]interface{}) []ValidationMessage {
  steps := runbook["steps"].([]interface{})
  var messages []ValidationMessage
  provisionIndices := make(map[string]int)

  // Find all provision steps
  for i, step := range steps {
    s := step.(map[string]interface{})
    id := s["id"].(string)
    if strings.HasPrefix(id, "provision-") {
      provisionIndices[id] = i
    }
  }

  // Check that each provision has a corresponding cleanup after it
  for provID, provIndex := range provisionIndices {
    cleanupID := strings.Replace(provID, "provision-", "cleanup-", 1)
    found := false
    for i := provIndex + 1; i < len(steps); i++ {
      s := steps[i].(map[string]interface{})
      if s["id"].(string) == cleanupID {
        found = true
        break
      }
    }
    if !found {
      messages = append(messages, ValidationMessage{
        File:     "<stdin>",
        Line:     provIndex + 5, // Approximate line number
        Severity: "error",
        Message:  fmt.Sprintf("Step '%s' requires a '%s' step after it", provID, cleanupID),
        Rule:     "compliance/cleanup-follows-provision",
      })
    }
  }

  return messages
}
\end{minted}

% ============================================================
\section{Error Reporting Format}
\label{sec:error-format}
% ============================================================

Validation messages MUST be structured JSON to enable tooling (IDEs, CI systems)
to parse and display them.

\paragraph{Field reference:}

\begin{description}
  \item[\texttt{file}] (string, required) --- File path. Use \texttt{"<stdin>"}
    if the runbook was piped in.
  \item[\texttt{line}] (integer, required) --- Line number (1-indexed).
  \item[\texttt{column}] (integer, optional) --- Column number (1-indexed).
  \item[\texttt{severity}] (string, required) --- \texttt{"error"} or \texttt{"warning"}.
    Errors block execution; warnings do not.
  \item[\texttt{message}] (string, required) --- Human-readable description.
    Should be actionable (tell the author what to fix).
  \item[\texttt{rule}] (string, required) --- Validator rule ID (e.g.,
    \texttt{"compliance/owner-required"}). Namespaced with Kit name.
\end{description}

\paragraph{Severity guidelines:}

\begin{itemize}
  \item Use \texttt{"error"} for violations that would cause incorrect behavior
    at runtime or violate hard requirements (e.g., missing required field,
    invariant violation).
  \item Use \texttt{"warning"} for stylistic issues, recommendations, or soft
    constraints (e.g., ``timeout exceeds recommended SLA'').
\end{itemize}

\paragraph{Message guidelines:}

\begin{itemize}
  \item Be specific: ``Step 'provision-db' requires a 'cleanup-db' step after it''
    is better than ``Invalid step order.''
  \item Be actionable: ``Add meta.roles.owner field'' is better than ``Owner missing.''
  \item Reference the rule ID in the message if it helps understanding.
\end{itemize}

% ============================================================
\section{Example: Validator in the Manifest}
\label{sec:validator-manifest}
% ============================================================

Validators are declared in the Kit manifest under the \texttt{validators} key:

\begin{minted}{yaml}
validators:
  - name: compliance/owner-required
    phase: parse                 # parse | plan
    description: "Every compliance runbook must declare an owner in meta.roles"
  - name: compliance/cleanup-follows-provision
    phase: plan
    description: "Every provision step must be followed by a cleanup step"
  - name: compliance/timeout-sla
    phase: parse
    description: "Warns if timeout exceeds recommended SLA"
\end{minted}

\paragraph{Validator invocation:}

Gert invokes validators by:

\begin{enumerate}
  \item Reading the Kit manifest to discover declared validators.
  \item Executing the Kit compiler with a \texttt{--validate} flag (or a
    separate validator binary if specified).
  \item Passing the runbook on stdin.
  \item Collecting validation messages from stdout.
  \item Aggregating messages from all validators and reporting them to the user.
\end{enumerate}

If your Kit bundles validators into the compiler binary, handle the
\texttt{--validate} subcommand:

\begin{verbatim}
./bin/gert-kit-compliance compile < runbook.yaml > lowered.yaml
./bin/gert-kit-compliance validate < runbook.yaml > errors.json
\end{verbatim}

If you prefer separate binaries per validator, structure them as:

\begin{verbatim}
./bin/validate-owner-required < runbook.yaml > errors.json
./bin/validate-cleanup-follows-provision < lowered.yaml > errors.json
\end{verbatim}
